problem:  
Let \( G \subset \mathrm{SO}(n+1) \) act by isometries on the round sphere \(S^n \subset \mathbb{R}^{n+1}\) with cohomogeneity one.  Suppose the slice representation of the principal isotropy \(H\) in \(G/H\) splits into precisely two non‑equivalent irreducible submodules \(V_1\oplus V_2\) of dimensions \(m_1,m_2>0\).  Denote by \(a_1(t),a_2(t)\) the metric radii associated with the two orbit directions along an arclength coordinate \(t\in(0,L)\) on the orbit space \([0,L]\).  Any \(G\)-invariant hypersurface \(M_\gamma \subset S^n\) arises as the preimage of a smooth curve \(\gamma(t)\subset(0,L)\).  Its mean curvature \(H_\gamma\) and scalar curvature \(S_\gamma\) satisfy an explicit second‑order ODE system depending on \((a_1(t),a_2(t))\) and their derivatives.

**Problem (Rigidity versus Multiplicity):**  
For cohomogeneity‑one \(G\)-actions on spheres with two irreducible isotropy summands as above, does the ODE system for \((H_\gamma=0, S_\gamma=\text{const})\) admit at most one non‑congruent \(G\)-invariant solution curve \(\gamma\)?  
Equivalently, is every \(G\)-invariant minimal hypersurface in the round \(S^n\) having constant scalar curvature uniquely determined (up to the group action and orientation), or can multiple distinct invariant examples coexist?

**Problem (Bifurcation aspect):**  
Fix such a \(G\)-action.  Take a one‑parameter family of \(G\)-invariant conformal metrics \(g_t=e^{2u_t}g_{\mathrm{round}}\) on \(S^n\), smooth in \(t\), with \(u_0=0\).  Does there exist \(t\neq0\) arbitrarily small for which a new branch of \(G\)-invariant minimal hypersurfaces with constant scalar curvature in \((S^n,g_t)\) bifurcates from a known solution at \(t=0\)? 

Why is it a "good" problem:  
This pair of questions pinpoints a sharply defined and plausible gap in the current theory of symmetric minimal submanifolds.  While Hsiang–Lawson theory reduces the minimality condition for a \(G\)-invariant hypersurface to an ODE system, the effect of imposing *constant scalar curvature* instead of general minimality has not been classified even for the simplest two‑summand actions (e.g., \(G=\mathrm{SO}(p+1)\times\mathrm{SO}(q+1)\) on \(S^{p+q+1}\)).  The rigidity/multiplicity problem asks whether this extra scalar‑curvature constraint isolates the familiar Clifford hypersurfaces or allows hidden families of new symmetric minimal embeddings—a deep geometric dichotomy that remains untested.  

The bifurcation version probes analytic stability under geometric perturbation of the ambient metric, revealing how conformal deformations might generate new invariant minimal hypersurfaces.  Both problems demand a synthesis of group‑theoretic classification, geometric analysis of nonlinear ODEs, and curvature identities.  They are expressible in elementary geometric terms yet could expose new rigidity phenomena linking symmetry, curvature, and variational geometry—an archetype of a “small‑entrance, deep‑depth” question that could open fresh lines of research.